\begin{lemma}
Let $D$ be a loopless $T_s$-free digraph on $N$ vertices.  If
\[
 \binom ns\left(\frac{\operatorname{fwi}_s(D)}{N^s}\right)^{\ell-1}<1,
\]
then the complete graph on $n$ vertices has an $\ell$-coloring with no
monochromatic $K_s$.
\end{lemma}
\begin{proof}
Choose independent uniform maps $\phi_c:[n]\to V(D)$ for
$c=1,\ldots,\ell-1$.  For $i<j$, give $ij$ the least color $c$ for which
$(\phi_c(i),\phi_c(j))$ is an arc, and give it color $\ell$ if no such color
exists.  Symmetrize the values on the reversed ordered pair; diagonal values
are irrelevant under \noderef{CompleteColoring}.

If an $s$-set were a clique in a color $c<\ell$, its vertices in increasing
order would map to a transitive tournament in $D$.  Looplessness ensures the
images are distinct, so this contradicts the hypothesis on $D$ and
\noderef{TransitiveTournament}.  For color $\ell$, a fixed increasing
$s$-tuple is monochromatic exactly when each of the $\ell-1$ independent maps
produces a forward-independent tuple.  Its probability is therefore
$(\operatorname{fwi}_s(D)/N^s)^{\ell-1}$, with repeated images already
included in the forward-tuple count.

By linearity of expectation, the expected number of color-$\ell$ cliques is
at most the displayed product.  The strict inequality makes this expectation
less than one, so some choice of the maps has no such clique.  The resulting
coloring has no monochromatic clique in any color.
\end{proof}
